\section{Data}

\subsection{Turnout Data}

We use state-level VEP (voting-eligible population) turnout from the
United States Elections Project~\citep{mcdonald2002turnout}, which
compiles official certified vote totals and Census-based VEP estimates.
VEP turnout is the preferred measure because it excludes ineligible voters
(non-citizens, felons in restricted states) from the denominator.

Table~\ref{tab:turnout-summary} reports state-level VEP turnout for the
five study states across four election cycles.

\begin{table}[h]
\centering\small
\caption{VEP turnout (\%) for study states, 2016--2022.
Source: United States Elections Project.}
\label{tab:turnout-summary}
\begin{tabular}{lccccc}
\toprule
State & 2016 & 2018 & 2020 & 2022 & Mean (cycles) \\
\midrule
NC & 64.5 & 49.7 & 74.5 & 48.3 & 59.3 \\
WI & 70.8 & 59.5 & 75.8 & 57.4 & 65.9 \\
TX & 51.6 & 46.3 & 60.0 & 38.8 & 49.2 \\
FL & 65.4 & 52.2 & 71.7 & 54.0 & 60.8 \\
PA & 69.2 & 52.8 & 72.8 & 47.8 & 60.7 \\
\midrule
Study state mean & 64.3 & 52.1 & 71.0 & 49.3 & 59.2 \\
\bottomrule
\end{tabular}
\end{table}

\paragraph{County-level approximation.}
For the district-level regression, we use county-level VEP estimates
derived from the state-level totals and county presidential vote shares
from VEST~\citep{vest2022}. This introduces measurement error in the
dependent variable; we treat the OLS coefficient as an attenuated estimate
of the true district-level effect.

\subsection{Compactness Data}

District-level Polsby-Popper scores are from the bisect pipeline's
compactness analysis (\texttt{bisect label-analyze --types compactness}).
For enacted plans: from the 2020 cycle enacted map.
For bisect plans: from the 2020 congressional configuration
(\texttt{bisect build official\_2020 --year 2020}).

National means: bisect plans 0.361, enacted gerrymanders 0.286.
